% ---
% Dedicatória
% ---
\begin{dedicatoria}
   \vspace*{\fill}
   \centering
   \noindent
   \textit{Este trabalho é dedicado a minha família e amigos\\
   que sempre me apoiaram nos momentos mais difíceis.} \vspace*{\fill}
\end{dedicatoria}
% ---

% ---
% Agradecimentos
% ---
\begin{agradecimentos}
Agradeço a minha esposa, Maria Clara, por todo apoio e pela base que me deu durante toda a minha caminhada. À minha mãe, agradeço por ter me mostrado a importância dos estudos e por ser um pilar em quem eu me tornei. 

Agradeço em especial ao meu orientador José Antonio Moreira Xexéo, pela paciência e pela orientação ao longo de toda a pesquisa, Agradeço também a Gabriela Moutinho de Souza Dias, Anderson F. Pereira dos Santos e Ronaldo Ribeiro Goldschmidt pelo apoio a este trabalho. 
Agradeço também ao corpo docente do Instituto Militar de Engenharia pelos ensinamentos que recebi ao longo da minha formação.


\end{agradecimentos}
% ---

% ---
% Epígrafe
% ---
\begin{epigrafe}
    \vspace*{\fill}
	\begin{flushright}
		\textit{``Não vos amoldeis às estruturas deste mundo, \\
		mas transformai-vos pela renovação da mente, \\
		a fim de distinguir qual é a vontade de Deus: \\
		o que é bom, o que Lhe é agradável, o que é perfeito.''\\
		(Bíblia Sagrada, Romanos 12, 2)}
	\end{flushright}
\end{epigrafe}
% ---

% ---
% RESUMOS
% ---

% resumo em português
\setlength{\absparsep}{18pt} % ajusta o espaçamento dos parágrafos do resumo
\begin{resumo}
\SingleSpacing
%iniciomudança
A padronização da família Ascon pelo National Institute of Standards and Technology (NIST), formalizada na publicação NIST SP 800-232, consolidou a Criptografia Leve (Lightweight Cryptography - LWC) como resposta à demanda por cifragem em dispositivos com restrições de processamento, memória e energia. Apesar da importância, por ser uma publicação recente, datada de agosto de 2025, a contagem de trabalhos com o objetivo de testar o Ascon em diferentes cenários ainda é baixa. Os trabalhos que medem até quantas rodadas esses algoritmos resistem, com aprendizado de máquina, em geral escolhem a diferença de entrada e testam a primitiva isolada, e não o que um adversário passivo observa no esquema AEAD completo. Esta dissertação mede o piso de rodadas do Ascon-AEAD128, do GIFT-COFB e do Schwaemm256-128, isto é, a maior quantidade de rodadas em que um classificador ainda separa o que o adversário observa de dados aleatórios, com o objetivo de indicar quanto cada algoritmo pode ser aliviado. O adversário é passivo e só vê criptogramas: duas mensagens seguidas do mesmo dispositivo, com nonces de contador consecutivos e textos em claro diferentes e desconhecidos. Uma primeira tentativa com criptogramas isolados de 64~KB e 641 atributos estatísticos ficou no acaso mesmo com o Ascon reduzido a uma rodada, já que a saída depende da chave e o sinal se dilui no criptograma. Com o XOR de pares de mensagens de 64 bytes de texto em inglês, 300 chaves com 100 pares cada, particionamento por chave (240 para treinamento e 60 para teste) e os classificadores Regressão Logística, Random Forest e XGBoost, os pisos foram de 3 das 12 rodadas no Ascon, 3 das 40 no GIFT-COFB e 2 das 11 no Schwaemm256-128, no pior caso testado, com o contador do dispositivo perto do zero. Na rodada seguinte ao piso nenhum detector superou o acaso, nem com busca de hiperparâmetros, redes neurais, dez vezes mais dados ou outras representações do par, e os braços de controle, com texto em claro aleatório, não tiveram detecção. Abaixo do piso as versões reduzidas são comprovadamente distinguíveis de dados aleatórios, e acima dele o resultado não é certificado de segurança.
%fimmudança

 \textbf{Palavras-chave}: \imprimirpalavraschave
\end{resumo}

% resumo em inglês
\begin{resumo}[Abstract]
 \begin{otherlanguage*}{english}
%  \linespread{1.3}
\SingleSpacing
%iniciomudança
The standardization of the Ascon family by the National Institute of Standards and Technology (NIST), formalized in NIST SP 800-232, consolidated Lightweight Cryptography (LWC) as a response to the demand for authenticated encryption on devices with processing, memory, and energy constraints. Despite its importance, few studies have tested Ascon across different scenarios so far, since the publication is recent, dated August 2025.
Studies that measure how many rounds these algorithms withstand, with machine learning or with statistical tests, usually choose the input difference and test the isolated primitive, not what a passive adversary observes in the full AEAD scheme. This dissertation measures the round floor of Ascon-AEAD128, GIFT-COFB, and Schwaemm256-128, that is, the largest number of rounds at which a classifier still separates what the adversary observes from random data, in order to indicate how much each algorithm can be lightened. The adversary is passive and sees only ciphertexts: two consecutive messages from the same device, with consecutive counter nonces and different, unknown plaintexts. A first attempt with isolated 64~KB ciphertexts and 641 statistical features stayed at chance even with Ascon reduced to one round, since the output depends on the key and the signal is diluted across the ciphertext. Using the XOR of pairs of 64-byte English text messages, 300 keys with 100 pairs each, key-based partitioning (240 keys for training and 60 for test), and Logistic Regression, Random Forest, and XGBoost classifiers, the floors were 3 of 12 rounds for Ascon, 3 of 40 for GIFT-COFB, and 2 of 11 for Schwaemm256-128, in the worst case tested, with the device counter near zero. In the round after the floor no detector exceeded chance, not even with hyperparameter search, neural networks, ten times more data, or other pair representations, and the control arms, with random plaintext, had no detections. Below the floor the reduced versions are demonstrably distinguishable from random data; above it, the result is not a security certificate.
%fimmudança
   \vspace{\onelineskip}

   \noindent
   \textbf{Keywords}: \imprimirkeywords
 \end{otherlanguage*}
\end{resumo}